\documentclass[11pt,a4paper]{article}

% Template visual Matriz Aprovação. Compile com XeLaTeX ou LuaLaTeX.
\usepackage{fontspec}
\usepackage{polyglossia}
\setmainlanguage{portuguese}
\IfFontExistsTF{Aptos}{
  \setmainfont{Aptos}
  \setsansfont{Aptos}
}{
  \IfFontExistsTF{Calibri}{
    \setmainfont{Calibri}
    \setsansfont{Calibri}
  }{
    \IfFontExistsTF{TeX Gyre Heros}{
      \setmainfont{TeX Gyre Heros}
      \setsansfont{TeX Gyre Heros}
    }{
      \setmainfont{Latin Modern Sans}
      \setsansfont{Latin Modern Sans}
    }
  }
}
\usepackage[a4paper,margin=2cm,headheight=16pt]{geometry}
\usepackage{graphicx}
\usepackage{xcolor}
\usepackage{tikz}
\usepackage{tcolorbox}
\usepackage{enumitem}
\usepackage{booktabs}
\usepackage{tabularx}
\usepackage{amssymb}
\usepackage{fancyhdr}
\usepackage{titlesec}
\usepackage{hyperref}

\definecolor{matrizlime}{HTML}{CBFF4D}
\definecolor{matrizink}{HTML}{101313}
\definecolor{matrizmuted}{HTML}{596363}
\definecolor{matrizline}{HTML}{DDE3DF}
\definecolor{matrizsoft}{HTML}{F3F7EC}

\hypersetup{colorlinks=true,linkcolor=matrizink,urlcolor=matrizink}
\pagestyle{fancy}
\fancyhf{}
\lhead{\textsf{\small MATRIZ APROVAÇÃO}}
\rhead{\textsf{\small APOSTILA}}
\cfoot{\textsf{\small \thepage}}
\renewcommand{\headrulewidth}{0.4pt}
\renewcommand{\headrule}{\hbox to\headwidth{\color{matrizline}\leaders\hrule height \headrulewidth\hfill}}

\titleformat{\section}{\Large\bfseries\color{matrizink}}{\thesection}{0.6em}{}
\titleformat{\subsection}{\large\bfseries\color{matrizink}}{\thesubsection}{0.6em}{}
\setlist{nosep,leftmargin=*}

\newtcolorbox{foco}{colback=matrizsoft,colframe=matrizlime,boxrule=1pt,arc=3pt,left=10pt,right=10pt,top=8pt,bottom=8pt}
\newtcolorbox{lei}{colback=matrizink,colframe=matrizink,coltext=white,boxrule=0pt,arc=3pt,left=10pt,right=10pt,top=8pt,bottom=8pt}
\newtcolorbox{alerta}{colback=white,colframe=matrizline,boxrule=0.6pt,arc=3pt,left=10pt,right=10pt,top=8pt,bottom=8pt}

% Logo usada na capa. Caminho relativo à pasta onde o .tex é compilado
% (materiais/<disciplina>/). Ajuste se a estrutura mudar.
\newcommand{\logomatriz}{../../public/logo.png}

\begin{document}

\begin{titlepage}
  \thispagestyle{empty}
  \begin{center}
    \includegraphics[width=0.55\linewidth]{\logomatriz}\par
  \end{center}
  \vspace{2.2cm}
  {\sffamily\fontsize{28}{32}\selectfont\bfseries Apostila — Contabilidade: Patrimônio e equação contábil\par}
  \vspace{0.4cm}
  {\sffamily\large Receita Federal, CGU, TCU/TCE e Banco Central\par}
  \vfill
  \begin{foco}
    \textsf{\textbf{Foco da apostila}}\\[3pt]
    todo lançamento deve preservar a equação \texttt{Ativo = Passivo Exigível + Patrimônio Líquido}.
  \end{foco}
  \vspace{0.7cm}
  {\sffamily\small Atualizado em 16 de agosto de 2026\quad ·\quad Cebraspe, FGV e FCC\par}
  \vspace{1cm}
  {\sffamily\small matrizaprova.com\par}
\end{titlepage}

\tableofcontents
\newpage

% Conteúdo gerado entra aqui.
\section*{O que cai}

Este é o fundamento para interpretar balanços e lançamentos. As bancas cobram:

\begin{enumerate}
  \item conceito de patrimônio;
  \item ativo, passivo e patrimônio líquido;
  \item equação fundamental da contabilidade;
  \item fatos contábeis permutativos, modificativos e mistos;
  \item origem e aplicação de recursos;
  \item diferença entre receita, despesa, custo e investimento.
\end{enumerate}

\textbf{Regra de prova:} todo lançamento deve preservar a equação \texttt{Ativo = Passivo Exigível + Patrimônio Líquido}.

\section*{Teoria essencial}

\subsection*{1. Patrimônio}

Patrimônio é o conjunto de bens, direitos e obrigações de uma entidade, avaliáveis em moeda.

\begin{itemize}
  \item \textbf{Bens:} elementos que podem ser utilizados ou controlados pela entidade e possuem valor econômico, como dinheiro, veículos, máquinas e estoques.
  \item \textbf{Direitos:} valores a receber de terceiros, como clientes, duplicatas a receber e empréstimos concedidos.
  \item \textbf{Obrigações:} dívidas e compromissos perante terceiros, como fornecedores, salários a pagar e empréstimos obtidos.
\end{itemize}

O patrimônio não se resume ao dinheiro disponível. Uma empresa pode ter pouco caixa e muitos direitos a receber, ou possuir bens relevantes financiados por obrigações.

\subsection*{2. Ativo}

Ativo é o conjunto de bens e direitos controlados pela entidade, resultantes de eventos passados e dos quais se espera que fluam benefícios econômicos futuros.

Exemplos:

\begin{itemize}
  \item caixa e bancos;
  \item contas a receber;
  \item estoques;
  \item máquinas e equipamentos;
  \item imóveis;
  \item aplicações financeiras;
  \item softwares e outros ativos intangíveis, quando reconhecidos conforme as normas aplicáveis.
\end{itemize}

O ativo representa, em linguagem didática, aplicações de recursos. Nem todo gasto vira ativo: é necessário que exista controle e expectativa de benefício econômico futuro, observados os critérios contábeis.

\subsection*{3. Passivo exigível}

Passivo é uma obrigação presente da entidade, derivada de eventos passados, cuja liquidação se espera que resulte na saída de recursos capazes de gerar benefícios econômicos.

Exemplos:

\begin{itemize}
  \item fornecedores a pagar;
  \item salários e tributos a recolher;
  \item empréstimos e financiamentos;
  \item provisões reconhecidas conforme os critérios aplicáveis.
\end{itemize}

O passivo exigível representa recursos de terceiros aplicados na entidade. A expressão "exigível" diferencia essas obrigações do patrimônio líquido.

\subsection*{4. Patrimônio líquido}

Patrimônio líquido é o valor residual dos ativos da entidade depois de deduzidos todos os seus passivos.

O patrimônio líquido pode incluir capital social, reservas, ajustes e resultados acumulados, conforme a estrutura societária e as normas aplicáveis.

Se o ativo é de R\$ 100.000 e o passivo exigível é de R\$ 70.000, o patrimônio líquido é de R\$ 30.000.

\subsection*{5. Equação fundamental}

A equação patrimonial é:

\texttt{Ativo = Passivo Exigível + Patrimônio Líquido}

Também pode ser escrita como:

\texttt{Patrimônio Líquido = Ativo - Passivo Exigível}

Exemplo:

\begin{tabularx}{\linewidth}{lX}
\toprule
 \textbf{Elemento} & \textbf{Valor} \\
\midrule
 Ativo & R\$ 250.000 \\
 Passivo exigível & R\$ 160.000 \\
 Patrimônio líquido & R\$ 90.000 \\
\bottomrule
\end{tabularx}


O balanço deve permanecer equilibrado após cada lançamento. Isso não significa que todos os lançamentos aumentem ou reduzam o patrimônio líquido; alguns apenas trocam elementos dentro do ativo ou do passivo.

\subsection*{6. Origens e aplicações}

Os recursos aplicados nos ativos têm origem em duas fontes principais:

\begin{itemize}
  \item capital próprio: patrimônio líquido;
  \item capital de terceiros: passivo exigível.
\end{itemize}

Se a empresa compra uma máquina à vista, aplica dinheiro em uma máquina. Se compra a prazo, a máquina é aplicação e a dívida com fornecedor é origem de terceiros.

\subsection*{7. Fatos contábeis}

Fatos contábeis são acontecimentos que provocam alteração no patrimônio ou podem ser registrados conforme a estrutura contábil.

\subsubsection*{7.1. Fato permutativo}

Altera a composição dos elementos patrimoniais, mas não modifica o patrimônio líquido.

Exemplo: compra de veículo à vista. O caixa diminui e o veículo aumenta pelo mesmo valor. O ativo total permanece igual.

\subsubsection*{7.2. Fato modificativo}

Altera o patrimônio líquido. Pode ser aumentativo, quando gera receita, ou diminutivo, quando gera despesa.

Exemplo aumentativo: recebimento de receita de serviço à vista. Aumentam caixa e patrimônio líquido.

Exemplo diminutivo: pagamento de despesa de aluguel. Reduzem caixa e patrimônio líquido.

\subsubsection*{7.3. Fato misto}

Combina permuta patrimonial com alteração do patrimônio líquido.

Exemplo: pagamento de uma obrigação com juros. Reduz caixa e passivo, mas os juros geram despesa e reduzem o patrimônio líquido.

\subsection*{8. Receita, despesa, custo e investimento}

\begin{itemize}
  \item \textbf{Receita:} aumento de benefícios econômicos que aumenta o patrimônio líquido, salvo aportes dos proprietários.
  \item \textbf{Despesa:} redução de benefícios econômicos que diminui o patrimônio líquido, salvo distribuições aos proprietários.
  \item \textbf{Custo:} gasto relacionado à produção de bens ou à prestação de serviços. Pode tornar-se despesa quando reconhecido no resultado.
  \item \textbf{Investimento:} aplicação de recursos com expectativa de benefício futuro, como aquisição de participação ou ativo.
  \item \textbf{Gasto:} sacrifício financeiro para obter bem ou serviço; pode virar custo, despesa ou investimento conforme sua finalidade.
\end{itemize}

O pagamento não define sozinho o momento da receita ou despesa. O regime de competência reconhece efeitos quando ocorrem, independentemente do pagamento ou recebimento, conforme as normas aplicáveis.

\subsection*{9. Balanço patrimonial}

O balanço patrimonial apresenta a posição patrimonial e financeira em uma data. De forma simplificada:

\begin{tabularx}{\linewidth}{lX}
\toprule
 \textbf{Ativo} & \textbf{Passivo + Patrimônio líquido} \\
\midrule
 bens e direitos & obrigações e recursos próprios \\
\bottomrule
\end{tabularx}


Ativo e passivo são classificados em circulante e não circulante conforme os critérios contábeis e o ciclo operacional. A classificação não deve ser feita apenas pela intuição: depende de prazo, natureza e condições da entidade.

\section*{Como a banca cobra}

\begin{itemize}
  \item Patrimônio líquido é residual: ativo menos passivo exigível.
  \item Compra à vista de bem é fato permutativo, em regra.
  \item Receita aumenta o patrimônio líquido; despesa reduz.
  \item Passivo representa obrigação, não simplesmente qualquer saída de caixa.
  \item Todo fato contábil deve preservar a equação patrimonial.
  \item Custo, despesa, gasto e investimento não são sinônimos.
  \item Pagamento pode ocorrer antes, depois ou no momento do reconhecimento contábil.
\end{itemize}

\section*{Exemplos resolvidos}

\subsection*{Exemplo 1}

Uma empresa possui ativos de R\$ 400.000 e obrigações de R\$ 250.000. Qual é o patrimônio líquido?

\textbf{R\$ 150.000.} Pela equação:

\texttt{PL = Ativo - Passivo = 400.000 - 250.000 = 150.000}.

\subsection*{Exemplo 2}

Uma empresa compra mercadorias por R\$ 20.000, pagando metade à vista e metade a prazo. Qual é o efeito inicial?

O estoque aumenta R\$ 20.000, o caixa diminui R\$ 10.000 e o passivo aumenta R\$ 10.000. O ativo total aumenta R\$ 10.000 e o passivo também aumenta R\$ 10.000, sem alteração imediata do patrimônio líquido.

\section*{Quadro-resumo}

\begin{tabularx}{\linewidth}{lX}
\toprule
 \textbf{Elemento} & \textbf{Conceito} \\
\midrule
 ativo & bens e direitos controlados \\
 passivo exigível & obrigações com terceiros \\
 patrimônio líquido & ativo menos passivo \\
 receita & aumenta o patrimônio líquido \\
 despesa & reduz o patrimônio líquido \\
 fato permutativo & troca elementos sem alterar o PL \\
 fato modificativo & altera o PL \\
 fato misto & combina troca e alteração do PL \\
\bottomrule
\end{tabularx}


\section*{Questões de fixação}

\subsection*{Questão 1 — (questão inédita)}

Se uma entidade possui ativo de R\$ 500.000 e passivo exigível de R\$ 320.000, seu patrimônio líquido é:

A) R\$ 180.000. B) R\$ 320.000. C) R\$ 500.000. D) R\$ 820.000. E) R\$ 160.000.

\begin{alerta}
\textbf{Gabarito: A.} \texttt{PL = Ativo - Passivo = 500.000 - 320.000 = 180.000}.
\end{alerta}

\subsection*{Questão 2 — (questão inédita)}

A compra de um veículo à vista, pelo valor de aquisição, é normalmente:

A) fato modificativo aumentativo. \\ B) fato modificativo diminutivo. \\ C) fato permutativo. \\ D) fato misto aumentativo. \\ E) fato sem registro contábil.

\begin{alerta}
\textbf{Gabarito: C.} Há redução de caixa e aumento de veículo, sem alteração imediata do patrimônio líquido.
\end{alerta}

\subsection*{Questão 3 — (questão inédita)}

Assinale o elemento que representa obrigação com terceiros:

A) estoque. B) caixa. C) clientes. D) fornecedores a pagar. E) capital social.

\begin{alerta}
\textbf{Gabarito: D.} Fornecedores a pagar representam passivo exigível.
\end{alerta}

\subsection*{Questão 4 — (questão inédita)}

Em regra, o reconhecimento de uma despesa provoca:

A) aumento do patrimônio líquido. \\ B) redução do patrimônio líquido. \\ C) aumento simultâneo do ativo e do patrimônio líquido. \\ D) redução obrigatória do passivo. \\ E) nenhum efeito patrimonial.

\begin{alerta}
\textbf{Gabarito: B.} A despesa reduz o resultado e, consequentemente, o patrimônio líquido, salvo situações específicas de reconhecimento e pagamento.
\end{alerta}

\subsection*{Questão 5 — (questão inédita)}

O patrimônio líquido é corretamente definido como:

A) total de bens, sem direitos e obrigações. \\ B) total das receitas recebidas. \\ C) valor residual dos ativos após dedução dos passivos. \\ D) total das dívidas com fornecedores. \\ E) soma do ativo e do passivo exigível.

\begin{alerta}
\textbf{Gabarito: C.} Essa é a definição contábil do patrimônio líquido.
\end{alerta}

\section*{Revisão final}

\begin{itemize}
  \item[$\square$] Sei diferenciar bens, direitos e obrigações.
  \item[$\square$] Memorizei a equação patrimonial.
  \item[$\square$] Sei calcular o patrimônio líquido.
  \item[$\square$] Diferencio capital próprio de capital de terceiros.
  \item[$\square$] Reconheço fatos permutativos, modificativos e mistos.
  \item[$\square$] Diferencio receita, despesa, custo, investimento e gasto.
  \item[$\square$] Sei que pagamento não define sozinho o reconhecimento contábil.
  \item[$\square$] Entendo a estrutura básica do balanço patrimonial.
\end{itemize}

\end{document}
